% !TEX root = ../main.tex
% ============================================================================
% Implementation & Experimental Setup
% ============================================================================
\chapter{Implementation and Experimental Setup}
\label{ch:impl}

Chapter~\ref{ch:method} defined the distillation objective mathematically. This
chapter describes how that objective is realised as a working training pipeline
(Section~\ref{sec:impl}) and the experimental protocol under which it is evaluated
(Section~\ref{sec:setup}). The implementation detail if of high importance 
because the distillation signal is subtle and easily corrupted: a mis-scaled,
mis-aligned, or numerically-unstable signal still produces a plausible-looking loss curve. 
So it is substantial to guarantee that only the intended
signal reaches the student.

\section{Implementation}
\label{sec:impl}

\subsection{Pipeline overview}
\label{ssec:impl_overview}
The distillation pipeline is built on top of the public HiVT codebase and adds four
stages: (i) an \emph{offline teacher-caching} pass that records the frozen
teacher's predictions for every training scene; (ii) a \emph{KD dataset/datamodule}
that joins those cached targets to the corresponding student inputs at load time;
(iii) a \emph{KD training module} that wraps an ordinary HiVT student and adds the
distillation term to its loss; and (iv) an \emph{evaluation and diagnostics} layer
that reports both geometric and calibration metrics. Stages are decoupled: because
the teacher is frozen, its forward passes are computed once and reused across every
student run and hyperparameter setting, so the marginal cost of an additional
student experiment is just the student's own training.

\subsection{Offline teacher caching}
\label{ssec:impl_cache}
The teacher (HiVT-128) is run once over the full training split and its predictions
are written to a single HDF5 store. For each scene, indexed by its sequence id, the
cache holds the focal agent's mixture parameters: the per-mode Laplace locations
($\loc^T_f$, shape $[F,H,2]$), the per-mode scales ($\scale^T_f$, $[F,H,2]$), and
the raw mode logits ($z^T_f$, $[F]$). Storing the parameters---rather than
re-running the teacher each step---makes the per-epoch teacher cost effectively
zero and decouples student training speed from teacher size, which is what allows
the $55\times$-smaller student to be trained at full speed against a large teacher.

A subtle correctness hazard lives in this stage. HiVT emits a trajectory for
\emph{every} actor in a scene in one batched forward pass, and the framework stacks
all scenes' actors into one node dimension; the focal agent must be recovered via
the batched focal-agent index. An early version of the cache indexed by scene
position instead, silently saving an arbitrary (mostly scene-0) actor's trajectory
under each sequence id---a defect that leaves training loss looking normal while
distilling against the wrong agent entirely. The cache is therefore written by
explicitly selecting the focal agent's node, in its rotated-local reference frame,
and the loss-formation step (Section~\ref{ssec:impl_loss}) re-asserts the
alignment. This is the practical face of the silent-failure problem flagged above.

Algorithm~\ref{alg:cache} states the caching pass; line~\ref{algline:focal} is the
focal-agent selection whose mis-implementation is the hazard just described.
Table~\ref{tab:cache_schema} documents the resulting on-disk layout. Two further
choices are worth noting. The cache stores the \emph{raw mode logits} $z^T_f$ rather
than the softmax weights $\pi^T_f$, deferring the normalisation 
\begin{equation}
  \pi^T_{f,k} = \frac{\exp(z^T_{f,k}/\tau)}{\sum_{j=1}^{F}\exp(z^T_{f,j}/\tau)}
  \label{eq:softmax-temp}
\end{equation}
to load time so the
same store can be read at any softmax temperature~$\tau$ (this work uses
$\tau{=}1$) and without committing a particular numerical form. And each write is guarded by 
a resume-safe existence
check (line~\ref{algline:resume}), so an interrupted pass---over a
$2\,{\times}\,10^5$-scene split this is a multi-hour job---resumes without
recomputing. The pass also emits two sidecar files: a \texttt{.manifest.json}
provenance record (teacher checkpoint, data root, split, local radius) and a
\texttt{.sanity\_stats.json} health summary (NaN, collapsed-scale, and
$\pi$-entropy checks), so a defective cache is caught before any student trains
against it.

\begin{algorithm}[htbp]
\caption{Offline teacher caching (run once over the training split)}
\label{alg:cache}
\begin{algorithmic}[1]
\Require frozen teacher $T$; training split $\mathcal{D}$; local radius $\rho$
  (matched to the student); output store $\mathcal{H}$
\For{each batch $b \in \mathcal{D}$}
  \State $(\mathrm{pred},\, z^T) \gets T(b)$
         \Comment{$\mathrm{pred}{:}\,[F,N,H,4]$ over all $N$ actors;\; $z^T{:}\,[N,F]$ logits}
  \State $\loc \gets \mathrm{pred}[\dots,{:}2]$;\quad
         $\scale \gets \mathrm{pred}[\dots,2{:}]$
         \Comment{$[F,N,H,2]$ each}
  \For{each scene $i$ in batch $b$}
    \State $\mathrm{sid} \gets \textsc{Normalize}(b.\mathrm{seq\_id}[i])$
    \If{$\mathrm{sid} \in \mathcal{H}$} \textbf{continue} \label{algline:resume}
         \Comment{resume-safe skip} \EndIf
    \State $n \gets b.\mathrm{agent\_index}[i]$ \label{algline:focal}
         \Comment{focal node in batched node space (PyG-offset)}
    \State $\mathcal{H}[\mathrm{sid}] \gets
            \big(\loc[:,n],\; \scale[:,n],\; z^T[n]\big)$
         \Comment{$[F,H,2],[F,H,2],[F]$, rotated-local frame}
  \EndFor
\EndFor
\State write \texttt{.manifest.json} (provenance) and \texttt{.sanity\_stats.json}
       (NaN / collapsed-scale / $\pi$-entropy checks)
\end{algorithmic}
\end{algorithm}

\begin{table}[htbp]
  \centering
  \caption[Teacher-cache layout.]{Per-scene contents of the teacher cache: one HDF5
    group per sequence id, three datasets, gzip-compressed. $F{=}6$ modes,
    $H{=}30$ steps, $2$ coordinates. Storing logits (not softmax weights) defers
    normalisation to load time.}
  \label{tab:cache_schema}
  \begin{tabular}{lllp{5.4cm}}
    \toprule
    Dataset & Shape & Dtype & Contents \\
    \midrule
    \texttt{loc}   & $[6,30,2]$ & float32 & per-mode Laplace means $\loc^T_f$ (focal agent, rotated-local frame) \\
    \texttt{scale} & $[6,30,2]$ & float32 & per-mode Laplace scales $\scale^T_f$ ($>0$) \\
    \texttt{pi}    & $[6]$      & float32 & raw mode logits $z^T_f$ (pre-softmax) \\
    \bottomrule
  \end{tabular}
\end{table}

\paragraph{Practical considerations.}
Two operational details proved to matter in practice. First, the cache was
initially written as one \texttt{.pt} file per scene; the resulting collection
of small files was wasteful both in disk footprint and in load overhead, and
consolidating it into a single HDF5 store reduced the on-disk size and replaced
thousands of small reads with indexed access into one handle. Second, the
offline-caching design trades teacher \emph{compute} for memory and I/O
pressure at load time: on a RAM-limited machine the operating system began
paging the working set to a mechanical disk, which collapsed training
throughput and, in one configuration, crashed the pipeline. Directing the
swap/page file to the SSD resolved this. The lesson is that the cost saved by
not re-running the teacher reappears as a storage-bandwidth requirement, so the
cache and any paging it induces should live on fast storage.

\subsection{The distillation loss module}
\label{ssec:impl_loss}
At train time, the KD module wraps a standard HiVT student and forms the combined
objective $\mathcal{L}=\lambda_{\text{task}}\,\mathcal{L}_{\mathrm{HiVT}} +
\lambda_{\text{kl}}\,\mathcal{L}_{\mathrm{KD}}$ (Eq.~\ref{eq:total}, with
$\lambda_{\text{task}}{=}1$ throughout and $\lambda_{\text{kl}}$ the swept
distillation weight). Three details make the term match its mathematical
definition:

\paragraph{Focal-agent slicing and an alignment assertion.}
The student's predictions are sliced to the focal agent using the same batched
index used to build the cache, and the resulting student tensor is checked against
the cached teacher tensor with an explicit shape assertion before the loss is
formed. Any batch/horizon/coordinate mismatch---the signature of a re-introduced
alignment bug---then fails loudly rather than training silently on garbage.

\paragraph{Mixture log-density in log-space.}
Both objectives evaluate the student mixture log-density
$\log p_S(\mathbf{y})=\log\sum_k \pi^S_k\,\Lap(\mathbf{y}\mid\loc^S_k,\scale^S_k)$.
This is computed entirely in log-space: per-element Laplace log-densities
$-\log(2b)-|y-\mu|/b$ are summed over horizon and coordinates, the student
log-mixture-weights (a log-softmax of the logits) are added, and the modes are
reduced with a numerically-stable \texttt{logsumexp}. This avoids underflow when
individual component densities are tiny---common when a teacher target lands far
from a student mode---which would otherwise produce \texttt{NaN} gradients.

\paragraph{Monte-Carlo sampling for v2.}
The mean-target loss (v1) scores the teacher locations $\loc^T_f$ directly under
$p_S$, probing the student at a single point per mode. The distribution-matching
loss (v2) must instead evaluate the expectation
$\mathbb{E}_{\mathbf{y}\sim\Lap(\loc^T_f,\scale^T_f)}[\log p_S(\mathbf{y})]$ of
Eq.~\ref{eq:v2}: the student's score $\log p_S$ averaged over \emph{everywhere} the
teacher mode believes the agent may go, weighted by how strongly it believes each
point. This average over a continuum of trajectories is an integral, and it is what
makes v2 reward a student that matches the teacher's whole spread rather than only
its centre. The integral has \emph{no closed form}, however: $p_S$ is itself a
log-sum-exp of Laplace densities, so a Laplace density times the log of a sum of
Laplaces cannot be integrated analytically. It is therefore estimated by Monte
Carlo, which replaces the integral by an average of its \emph{integrand}---the
quantity inside the integral, here the student score $\log p_S$---evaluated at
$S{=}8$ samples drawn from the teacher mode,
\begin{equation}
  \label{eq:mc}
  \mathbb{E}_{\mathbf{y}\sim\Lap(\loc^T_f,\scale^T_f)}[\log p_S(\mathbf{y})]
    \;\approx\; \frac{1}{S}\sum_{s=1}^{S}\log p_S\!\left(\mathbf{y}^{(s)}\right),
  \qquad
  \mathbf{y}^{(s)}\sim\Lap(\loc^T_f,\scale^T_f).
\end{equation}
A subtlety distinguishes this from merely relocating the target, and it is the
reason v2 behaves differently from v1 at all. Averaging the \emph{samples}
$\mathbf{y}^{(s)}$ would simply recover the mean $\loc^T_f$---they are drawn
symmetrically about it---collapsing v2 back to v1. What Eq.~\eqref{eq:mc} averages
is instead the \emph{integrand} $\log p_S(\mathbf{y}^{(s)})$, the student's score
\emph{at} each sample. Because $\log p_S$ is nonlinear (a log-sum-exp), the mean of
the scores is not the score of the mean,
\[
  \tfrac1S\textstyle\sum_s \log p_S\!\left(\mathbf{y}^{(s)}\right)
  \;\neq\; \log p_S\!\left(\loc^T_f\right),
\]
and the gap between the two carries exactly the information about the teacher's
spread $\scale^T_f$ that v1 throws away. Concretely, when a teacher mode is
uncertain (large $\scale^T_f$) its samples are thrown wide, so a student that has
collapsed its density onto a narrow spike scores highly only near $\loc^T_f$ and
poorly on the outlying samples, dragging the average down; only a student whose
density covers the teacher's spread keeps every sample's score high. The estimator
thus penalises precisely the over-confident scale shrinkage that v1 cannot detect,
and in the zero-spread limit $\scale^T_f\!\to\!0$ every sample coincides with
$\loc^T_f$ and v2 reduces exactly to v1.

Samples are drawn by the inverse-CDF (reparameterised) method rather than a
black-box sampler: drawing $u\sim\mathcal{U}(-\tfrac12,\tfrac12)$ independently per
element and applying the inverse Laplace CDF gives an exact draw,
\begin{equation}
  \label{eq:invcdf}
  \mathbf{y}^{(s)} \;=\; \loc^T_f \,-\, \scale^T_f \odot
    \operatorname{sgn}(u)\odot\log\!\big(1-2\lvert u\rvert\big).
\end{equation}
Inverse-transform sampling is what makes this an \emph{exact} Laplace draw. A
distribution's CDF maps a value to the probability mass below it; feeding a uniform
$u$ \emph{backwards} through the inverse CDF returns the value at that percentile,
and because $u$ is uniform over percentiles the returned values follow the target
distribution---percentiles near the median, which lie geometrically close to
$\loc^T_f$, are requested often and the tails rarely, reproducing the Laplace shape
with no rejection step. The Laplace is one of the few distributions whose inverse
CDF is closed-form, which collapses the whole procedure into the single expression
above: $u\!+\!\tfrac12$ is the requested percentile, $\operatorname{sgn}(u)$ selects
the side of the mode, $\log(1-2\lvert u\rvert)$ sets the distance out (near $0$ at
the median, diverging towards the tails), and $\scale^T_f$ scales that distance so
the teacher's own uncertainty governs how far the samples spread. The percentiles
are not a fixed grid but eight independent random draws per mode, re-rolled each
step---the ``Monte-Carlo'' in the name---so the estimate fluctuates from step to
step with the variance noted below.

Writing the sample as a deterministic, differentiable transform of the parameters
and an external noise source $u$ is the reparameterisation trick. It is worth being
precise about \emph{why} it matters here, because the usual motivation does not
apply directly: the expectation is taken over the \emph{teacher} distribution, which
is frozen and cached, while the student parameters enter only through the integrand
$\log p_S$. The student gradient $\nabla_S\frac1S\sum_s\log p_S(\mathbf{y}^{(s)})$ is
therefore exact for any fixed samples, therefore no score-function correction through the
sampling distribution is needed, so inverse-CDF sampling is used because it is
exact, cheap, and fully vectorised, and because the reparameterised form keeps the
estimator differentiable should the same code ever be reused with an unfrozen or
learned teacher. The estimator~\eqref{eq:mc} is unbiased and its variance falls as
$1/S$; $S{=}8$ gives a stable gradient at an inner cost of $O(S K_S H)$ per teacher
mode. Each mode's contribution is then weighted by its softmax confidence $\pi^T_f$,
realising the confidence-weighted soft-target set of
Section~\ref{ssec:method_pi}.

In code this is implemented batched rather than as the per-mode loop of
Algorithm~\ref{alg:v2}: the teacher samples are flattened across the
sample and mode axes into a single batch of $S\,K_T$ soft targets, scored in one
call to the \emph{same} mixture-density routine that v1 uses, and only then
reshaped and averaged back over the $S$ samples. Figure~\ref{fig:v2_pipeline}
traces this dataflow end to end with the tensor shape at each stage; it makes
explicit that v2 is v1 wrapped in a sample--flatten--average envelope, and that the
two variants share the single scoring kernel highlighted in the figure.

\begin{figure}[htbp]
  \centering
  \begin{tikzpicture}[
      node distance=9mm and 0mm,
      box/.style={draw, rounded corners, align=center, inner sep=5pt,
                  minimum width=46mm, font=\small},
      kernel/.style={box, draw=black, fill=black!6, very thick},
      slbl/.style={font=\footnotesize, align=left, anchor=west},
      note/.style={font=\footnotesize\itshape, align=left, anchor=west},
      arr/.style={-{Stealth[length=2.4mm]}, thick},
    ]
    \node[box] (teacher)
      {teacher modes $\loc^T,\scale^T$\\[2pt]\texttt{[F,\,B,\,H,\,D]}};
    \node[box, below=of teacher] (samp)
      {samples $\mathbf{y}\sim\Lap(\loc^T,\scale^T)$\\[2pt]\texttt{[S,\,F,\,B,\,H,\,D]}};
    \node[box, below=of samp] (flat)
      {flattened soft targets\\[2pt]\texttt{[S*F,\,B,\,H,\,D]}};
    \node[kernel, below=of flat] (score)
      {$\log p_S(\cdot)$ \;(v1 mixture scorer)\\[2pt]\texttt{[S*F,\,B]}};
    \node[box, below=of score] (avg)
      {$\overline{\log p_S}$ over $S$\\[2pt]\texttt{[F,\,B]}};
    \node[box, below=of avg] (out)
      {$\mathcal{L}_{\mathrm{v2}}$\;\;(scalar)};

    \draw[arr] (teacher) -- node[slbl, midway, xshift=4mm]
      {(1) draw $S$ samples per\\[-1pt]\hphantom{(1) }mode (Eq.~\ref{eq:invcdf})} (samp);
    \draw[arr] (samp) -- node[slbl, midway, xshift=4mm]
      {(2) flatten sample$\times$mode} (flat);
    \draw[arr] (flat) -- node[slbl, midway, xshift=4mm]
      {(3) score under student\\[-1pt]\hphantom{(3) }mixture $p_S$} (score);
    \draw[arr] (score) -- node[slbl, midway, xshift=4mm]
      {(4) average over the $S$\\[-1pt]\hphantom{(4) }samples (Eq.~\ref{eq:mc})} (avg);
    \draw[arr] (avg) -- node[slbl, midway, xshift=4mm]
      {(5) $\times\,\pi^T_f$, sum over $F$,\\[-1pt]\hphantom{(5) }$\div 2H$, mean over $B$, $\times\,\lambda$} (out);

    \node[note, right=18mm of avg]
      {identical shape and\\scorer as v1 from here};
  \end{tikzpicture}
  \caption[Batched dataflow of the v2 distillation loss.]{Batched tensor flow of
    the distribution-matching loss $\mathcal{L}_{\mathrm{v2}}$, with the shape at
    each stage. $F$ teacher modes, $B$ focal agents in the batch, $H{=}30$ horizon
    steps, $D{=}2$ coordinates, $S{=}8$ Monte-Carlo samples. The shaded box is the
    student mixture scorer $\log p_S$, the \emph{same} kernel evaluated by v1
    (Algorithm~\ref{alg:v1}); v2 only adds the sample (1) and flatten (2) before it
    and the average (4) after it. From stage (4) onward the computation is identical
    in shape and code to v1, which scores the teacher means directly. Step (5)
    folds in the teacher confidences and the per-element normalisation of
    Algorithm~\ref{alg:v2}.}
  \label{fig:v2_pipeline}
\end{figure}

\subsection{Evaluation and silent-failure diagnostics}
\label{ssec:impl_diag}
Because a wrong distillation signal need not perturb the loss curve, the pipeline is
instrumented with diagnostics designed to expose exactly the failure modes the
method is sensitive to. Alongside the geometric metrics, evaluation reports the mean
predicted scale ($b$-scale, to detect over-confident shrinkage), reliability and
coverage curves (to detect mis-calibration), mode-weight entropy (to distinguish a
scale pathology from a collapse of the mixture weights), and an offline
mode-diversity measure (to detect a student mimicking only the teacher's dominant
mode). A separate, training-free diagnostic verifies the premise of the whole method
by computing the optimal per-scene teacher--student mode assignment
(Section~\ref{ssec:method_perm}); its output is the assignment matrix of
Figure~\ref{fig:perm}. These diagnostics are what turn the theoretical predictions
of Chapter~\ref{ch:method} into falsifiable checks, and their results are reported
throughout Chapter~\ref{ch:results}.

\section{Experimental Setup}
\label{sec:setup}

\subsection{Dataset}
\label{ssec:setup_data}
All experiments use the Argoverse~1 motion-forecasting
benchmark~\cite{chang2019argoverse}, the dataset HiVT was designed for. Each
sequence provides the tracked trajectories of all agents in an urban driving scene
together with an HD map, sampled at $10$~Hz. The task is to observe the first
$2$~seconds ($20$ steps) and predict the next $3$~seconds ($30$ steps, so $H=30$)
for a single designated focal agent. The benchmark provides $205{,}942$ training
and $39{,}472$ validation sequences. Models are trained on the training split and
all reported numbers are computed on the validation split, on which ground truth is
available; the held-out test split is evaluated only through the official
leaderboard and is not used here.

\subsection{Metrics}
\label{ssec:setup_metrics}
Following the benchmark, the model emits $K=6$ trajectory modes per agent.
Evaluation spans two complementary groups.

\paragraph{Geometric (best-of-$K$) accuracy.}
These measure whether \emph{any} predicted mode is close to the ground truth.
\begin{itemize}[leftmargin=1.4em]
  \item \textbf{minADE}$_K$ --- the average displacement error (mean over the $H$
        steps) of the best of the $K$ modes.
  \item \textbf{minFDE}$_K$ --- the displacement error at the final predicted step
        of the best mode; the primary geometric metric in this work.
  \item \textbf{MR} (miss rate) --- the fraction of scenes whose best mode's final
        point exceeds $2$~m from the ground truth.
  \item \textbf{brier-minFDE} --- the official Argoverse ranking metric, minFDE
        augmented by a Brier penalty $(1-\pi_{k^\star})^2$ that rewards assigning
        high probability to the best mode, thereby coupling geometry to the mode
        weights.
\end{itemize}

\paragraph{Calibration (full-distribution) quality.}
Best-of-$K$ metrics ignore the predicted \emph{uncertainty}; the following probe
whether the full mixture is honest, which matters because a downstream planner
consumes the distribution, not just the best mode.
\begin{itemize}[leftmargin=1.4em]
  \item \textbf{mixNLL} --- the negative log-likelihood of the ground truth under
        the full predicted mixture~\eqref{eq:mixture}, defined in~\eqref{eq:mixnll};
        the single scalar that rewards both accuracy and well-calibrated spread, and
        the calibration metric logged by the training monitor. It must not be
        confused with the best-mode regression NLL of~\eqref{eq:reg}, which scores
        only the winning mode.
  \item \textbf{calibration error} --- the discrepancy between nominal and empirical
        coverage of the predicted intervals (deviation of a reliability curve from
        the diagonal).
  \item \textbf{coverage at $p$} (e.g.\ cov@p90) --- the empirical fraction of
        ground-truth points falling within the predicted $p$ central interval;
        well calibrated when it matches the nominal $p$.
  \item \textbf{$b$-scale} --- the mean predicted Laplace scale, a direct measure of
        sharpness used to diagnose over-confidence (scale shrinkage).
  \item \textbf{mode-weight entropy} ($\pi$-entropy) --- the entropy of
        $\{\pi_k\}$, used to distinguish scale pathologies from collapse of the
        mode weights.
\end{itemize}

Table~\ref{tab:metrics} collects every metric used in this work with its formula and
its direction of improvement, and serves as the reference index for the symbols that
recur throughout Chapter~\ref{ch:results}.

\begin{table}[htbp]
  \centering
  \caption{Index of evaluation metrics. $K{=}6$ modes over $H$ prediction steps;
    $\mathbf{y}^\star$ is the ground-truth future,
    $k^\star=\argmin_k\sum_t\lVert\loc_{k,t}-\mathbf{y}^\star_t\rVert_2$ the best mode
    and $\pi_{k^\star}$ its weight. Geometric metrics use the single best mode;
    calibration metrics use the full mixture~\eqref{eq:mixture}. ``Better'' gives the
    desired direction ($\downarrow$ lower, $\to p$ should match the nominal level,
    ``diag.'' a diagnostic read alongside the others).}
  \label{tab:metrics}
  \small
  \begin{tabular}{llp{6.6cm}c}
    \toprule
    Metric & Symbol & Definition & Better \\
    \midrule
    \multicolumn{4}{l}{\emph{Geometric (best-of-$K$) accuracy}} \\
    minADE$_K$    & --- & $\frac{1}{H}\sum_{t=1}^{H}\lVert\loc_{k^\star,t}-\mathbf{y}^\star_t\rVert_2$, mean displacement of the best mode & $\downarrow$ \\
    minFDE$_K$    & --- & $\lVert\loc_{k^\star,H}-\mathbf{y}^\star_H\rVert_2$, final-step error of the best mode (primary geometric metric) & $\downarrow$ \\
    MR            & --- & fraction of scenes with $\text{minFDE}>2\,\mathrm{m}$ & $\downarrow$ \\
    brier-minFDE  & --- & $\text{minFDE}+(1-\pi_{k^\star})^2$, the official Argoverse ranking metric & $\downarrow$ \\
    \addlinespace
    \multicolumn{4}{l}{\emph{Calibration (full-distribution) quality}} \\
    mixNLL        & $\mathcal{L}_{\mathrm{mix}}$ & $-\log\sum_{k}\pi_k\,\Lap(\mathbf{y}^\star\mid\loc_k,\scale_k)$, Eq.~\eqref{eq:mixnll} & $\downarrow$ \\
    calib.\ error & --- & deviation of the reliability curve from the diagonal (nominal vs.\ empirical coverage) & $\downarrow$ \\
    cov@$p$       & --- & empirical fraction of $\mathbf{y}^\star$ inside the predicted $p$ central interval & $\to p$ \\
    $b$-scale     & $\bar{b}$ & mean predicted Laplace scale (sharpness; small $\Rightarrow$ over-confident) & diag. \\
    $\pi$-entropy & --- & $-\sum_{k}\pi_k\log\pi_k$, entropy of the mode weights & diag. \\
    \bottomrule
  \end{tabular}
\end{table}

\subsection{Models and Training Protocol}
\label{ssec:setup_models}
The capacity axis is HiVT's embedding width $d$ (Section~\ref{ssec:sota_hivt}),
which sets the dimension of essentially every weight matrix and hence the whole
parameter count. Four widths are studied, summarised in Table~\ref{tab:models}:
$d\in\{128,64\}$ from the authors' released checkpoints serve as the large
references (and $d{=}128$ as the distillation teacher), while $d\in\{32,16\}$ are
trained in this work as the student candidates.

\begin{table}[htbp]
  \centering
  \caption{HiVT configurations studied. Widths $128$ and $64$ use the authors'
    released checkpoints; $32$ and $16$ are trained in this work. Parameter counts
    are exact (architecture-determined).}
  \label{tab:models}
  \begin{tabular}{lrrr}
    \toprule
    Model       & Width $d$ & Parameters & vs.\ teacher \\
    \midrule
    HiVT-128 (teacher) & 128 & 2{,}559{,}993 & $100\%$ (1$\times$) \\
    HiVT-64            & 64  & 653{,}369     & $25.5\%$ (4$\times$ smaller) \\
    HiVT-32 (student)  & 32  & 170{,}073     & $6.6\%$ (15$\times$ smaller) \\
    HiVT-16 (student)  & 16  & 45{,}929      & $1.8\%$ (55$\times$ smaller) \\
    \bottomrule
  \end{tabular}
\end{table}

\paragraph{Attention heads.}
HiVT uses $8$ attention heads at every width by default, so the per-head
dimension $d/\text{heads}$ shrinks with width: $128/8{=}16$ at the teacher,
$64/8{=}8$ at HiVT-64, and only $32/8{=}4$ at HiVT-32. HiVT-32 keeps the default
$8$ heads (head dimension $4$), so the comparison against the teacher varies width
alone; halving this to $2$ heads (head dimension $16$, matching the teacher's
per-head width) was also examined and is discussed in
Chapter~\ref{ch:results}. At $d{=}16$, however, the default $8$ heads give a head
dimension of only $2$, which is too small for the attention to be expressive:
training stalls at a degenerate optimum (minADE plateaus near $3.3$, against
roughly $1$ for a working configuration). HiVT-16 therefore uses $2$ heads (head
dimension $8$), and this is the configuration behind every reported HiVT-16
result.

\paragraph{Optimisation.}
All models are trained with AdamW (Section~\ref{ssec:bg_optimizers}) under a
cosine-annealing schedule (Section~\ref{ssec:bg_schedulers}) for $64$ epochs at
batch size $128$. The learning rate is the main width-dependent hyperparameter:
HiVT-32 trains best at $\eta=3\times10^{-3}$, whereas HiVT-16 underfits at that
rate and uses $\eta=1\times10^{-3}$ (both selected by a learning-rate sweep
reported in Chapter~\ref{ch:results}). Training is performed on a single NVIDIA
RTX~5060; a full $64$-epoch run on the complete Argoverse~1 training split takes
roughly $19$~hours for HiVT-32 (and less for the smaller HiVT-16).

\subsection{Distillation Protocol}
\label{ssec:setup_kd}
The teacher is always HiVT-128, with its predictions cached offline
(Section~\ref{ssec:impl_cache}). For each student width a \emph{from-scratch}
baseline ($\lambda=0$, no distillation) is trained under an identical recipe, so
that the effect of distillation is measured by an A/B comparison in which only the
distillation term changes. Both distillation variants
(Section~\ref{ssec:method_v1},~\ref{ssec:method_v2}) are evaluated, with the
distribution-matching loss using $S=8$ Monte-Carlo samples per teacher mode, and
the weight $\lambda$ swept per variant and width (recalling that $\lambda$ is not
comparable across variants). To contain compute, hyperparameters are first ranked
on a $25\%$-data, short-schedule \emph{proxy} and the selected configuration is
then confirmed on the full dataset at the full $64$-epoch schedule; both the proxy
and the confirmation are reported in Chapter~\ref{ch:results}.
